\vsssub
\subsubsection{~分布式内存概念} \label{sec:distr}
\vsssub

前一段描述的一般网格结构同时用于模式的共享内存版本和分布式内存版本，
只存在少量差异。对于模式的分布式内存版本，并非所有数据都保存在每个
处理器上。相反，每个谱只保存在单个处理器上。存储网格上的谱以固定步长
分布到可用处理器上。由于在给定处理器上只本地存储部分谱，需要区分上文
的全局海点计数器 {\F isea} 和本地海点计数器 {\F jsea}。如果计算中实际
使用的处理器数量为 {\F naproc}，而 {\F iaproc} 是从 1 到 {\F naproc} 的
处理器编号，则这些参数的关系为

\vspace{\baselineskip}
\centerline{\F isea = iaproc + (jsea-1) naproc ,}
\centerline{\F jsea = 1 + (isea-1) / naproc ,}
\centerline{\F iaproc = 1 + mod(isea-1,naproc) .}
\vspace{\baselineskip}

\noindent
在模式 3.10 版中，引入了进一步细化。实际处理器数量 {\F naproc} 可以小于
程序使用的处理器总数（{\F ntproc}）。满足 {\F naproc} $<$ {\F iaproc}
$\le$ {\F ntproc} 的处理器仅保留用于输出处理。

采用这种数据分布时，源项和谱内传播可以在每个给定处理器上计算，而不需要
处理器之间通信。不过，对于空间传播，需要进行数据转置：必须把所有空间
网格点的谱分量 {\F(ith,ik)} 汇集到单个处理器上。传播执行后，修改后的
数据必须再分散回其“归属”处理器。各个谱分量分配给特定处理器的方式，使得
每个处理器需要执行的部分传播步数大致相同。这使得良好的负载均衡成为可能。
实际算法可见子程序 {\F w3init}（{\file w3initmd.ftn}）的 4.d 节。

汇集操作的数据转置使用 Message Passing Interface (MPI) 标准
\citep[例如][]{bk:GLS97} 分两步实现。首先，对于给定谱 bin {\F(ith,ik)}，
每个空间网格点的值在单个目标处理器上汇集到一维数组 {\F store(isea)} 中，
随后该数组被转换为谱分量的完整二维场。传播执行后，分散操作的转置会使用
同一个一维数组 {\F store} 反向执行该过程。虽然把空间传播分配到各处理器的
算法保证了全局（每时间步）负载均衡，但它不能保证通信同步，因为每个处理器
上的每项计算并不一定耗费相同工作量。为避免由此产生负载不均衡，使用了
非阻塞通信。此外，一维数组 {\F store(isea)} 被替换为 {\F store(isea,ibuf)}，
其中新增数组维度提供主动管理的缓冲空间（见 {\file w3wavemd.ftn} 中的
{\F w3gath} 和 {\F w3scat}）。这些缓冲区允许在通信期间可能出现的空闲时钟
周期用于计算；如果硬件具备能力，也会把通信隐藏在计算之后。为避免 FORTRAN
与 MPI 之间的不兼容问题，使用了独立的汇集和分散数据数组。带缓冲的数据
转置以图形方式示于图~\ref{fig:transpose}。更多细节见 \cite{tol:PACO02}。

% fig:transpose

\setlength{\unitlength}{0.1mm}
\begin{figure}

\begin{picture}(1370,650)(0,-650)

\put(   0, -650){\framebox(1370,650)[c]{}}

\put( 270, -160){\vector(3,-4){150}}
\put( 270, -260){\vector(3,-2){150}}
\put( 270, -360){\vector(1, 0){150}}
\put( 270, -460){\vector(3, 2){150}}
\put( 270, -560){\vector(3, 4){150}}

\put( 820, -360){\vector(1,0){120}}

\put( 420, -380){\vector(-3, 4){150}}
\put( 420, -380){\vector(-3, 2){150}}
\put( 420, -380){\vector(-1, 0){150}}
\put( 420, -380){\vector(-3,-2){150}}
\put( 420, -380){\vector(-3,-4){150}}

\put( 940, -380){\vector(-1,0){120}}

\put(  10,  -80){\makebox(310,80)[c]{带有原生数据的}}
\put(  10, -110){\makebox(310,80)[c]{处理器}}
\put(  60, -210){\framebox(210,80)[c]{{\F 1}}}
\put(  60, -310){\framebox(210,80)[c]{{\F 2}}}
\put(  60, -410){\framebox(210,80)[c]{{\F 3}}}
\put(  60, -510){\framebox(210,80)[c]{{\F \ldots}}}
\put(  60, -610){\framebox(210,80)[c]{{\F naproc}}}

\put( 370,  -80){\makebox(500,80)[c]{含单个谱分量的}}
\put( 370, -120){\makebox(500,80)[c]{一维全网格数组}}
\put( 420, -400){\framebox(400,60)[c]{活动}}
\put( 420, -460){\dashbox{5}(400,240)[c]{}}
\put( 430, -280){\makebox(380,60)[l]{缓冲空间}}

\put( 970,  -80){\makebox(270,80)[c]{对应的}}
\put( 970, -120){\makebox(270,80)[c]{二维数组}}
\put( 940, -520){\framebox(330,300)[c]{传播}}

\put( 370, -570){\dashbox{15}(950,400)[c]{}}
\put( 400, -580){\makebox(270,80)[l]{\small 在目标处理器上}}

\end{picture}

\caption{分布式内存模式版本中的数据转置。首先，在汇集操作中，数据在图中
从左向右移动。计算完成后，数据在分散操作中从右向左移动。}
\label{fig:transpose}

\botline
\end{figure}


原则上，只有存储数组 {\F a(ith,ik,jsea)} 受数据分布影响。输入场、映射和
平均波浪参数输出场原则上在每个网格点都保留完整分辨率。计算的每个阶段，
每个处理器上都可获得完整映射。输入和输出场通常只在步长 {\F naproc} 处
包含相关数据。

分布式内存还需要修改 I/O。输入文件由每个独立处理器完整读取。文件输出
类型由 I/O 类型指示符 {\F iostyp} 决定。

\begin{center} \begin{tabular}{cl}
{\F iostyp} & 表示 \\ \hline
  0 & 每个独立进程写出重启动文件。 \\
  1 & 每个文件由分配到的进程写出。 \\
  2 & 每个文件由单个专用输出进程写出。 \\
  3 & 每类输出使用专用输出进程。
\end{tabular} \end{center}

\noindent
注意，重启动文件是直接访问文件，因此每个处理器可以高效地只汇集本地存储
的谱，而不需要通读整个文件。重启动文件可以由每个独立进程直接写出，也可
通过专用处理器汇总所有数据。第一种方法需要并行文件系统，第二种方法通常
适用。

当前数据分布算法的选择有几个原因。首先，它针对源项的（动态）积分、冰覆盖
网格点的排除以及谱内传播，提供自动且高效的负载均衡。其次，与更传统的区域
分解不同，这种通信按定义独立于数值传播格式。在后一种情况下，只需把所谓
边界数据“晕区”转换到相邻的网格点“块”中。晕区大小取决于所选传播格式。
当前数据分布方案的主要缺点是，每个时间步需要通信的数据量远大于更传统的
区域分解，特别是在使用处理器数量相对较少时。在测试 \ws\ 分布式内存版本的
IBM RS6000 SP 上，相对较大的通信量并未构成总计算时间的重要部分，并且模式
在多达 O(100) 个处理器时表现出优良的扩展行为 \citep{tol:PACO02}。

最近，又开发了混合并行化技术：结合粗尺度区域分解和本地数据转置，并使用
{\file ww3\_multi} 中已有的方法。为支持这种做法，模式 4.10 版引入了
{\file ww3\_gspl(.sh)} 工具。虽然这种方法在 {\file ww3\_multi} 初始化阶段的
模式内存占用方面仍需进一步工作，但用该方法获得的初步扩展结果令人鼓舞
\cite[见][]{tol:MMAB13b}。
